\graphicspath{{abstracts/praesentationen/02-migration-sammlungsdatenbank/img/}{img/}}

\begin{beitrag}{Vom Excel-Sheet zur Triple-Store-Datenbank: Migration einer historischen Sammlungsdatenbank}{%
Clara Castorp\,\orcidlink{0000-0003-4567-8901} \\
\small \href{mailto:clara.castorp@example.org}{clara.castorp@example.org} \\
\small Universität Heidelberg \\[0.5cm]
Daniel Drechsler\,\orcidlink{0000-0001-5678-9012} \\
\small \href{mailto:daniel.drechsler@example.org}{daniel.drechsler@example.org} \\
\small Humboldt-Universität zu Berlin
}

\section{Ausgangslage}
Sammlungsdatenbanken in kleineren Forschungsinstituten sind häufig historisch gewachsen: Eine erste Aufnahme in einer relationalen Datenbank, eine zweite Erschließung mit Hilfe von Tabellenkalkulationen, eine spätere Anreicherung in einem Wiki -- am Ende liegen die Bestände in drei oder mehr unverbundenen Systemen vor. Im Fall der von uns betreuten Sammlung historischer Druckschriften des 18.\,Jahrhunderts mussten 12\,400 Datensätze aus einem Excel-Bestand, 3\,200 OCR-erfasste Inhaltsregister und eine externe Provenienz-Datenbank in eine konsolidierte, WissKI-basierte Triple-Store-Lösung überführt werden. Der vorliegende Beitrag berichtet über die methodischen Schritte, die Werkzeugkette und die Lehren aus diesem Migrationsprojekt.

\section{Methodisches Vorgehen}
Die Migration erfolgte in vier Etappen. Im ersten Schritt wurden die Quelldaten profiliert: Datentypen, Wertebereiche, implizite Schlüsselbeziehungen und Inkonsistenzen wurden mit \texttt{OpenRefine} und einer projektspezifischen Skriptsammlung erfasst. Bereits in dieser Phase wurde deutlich, dass ein erheblicher Anteil der Datenqualitätsprobleme auf nicht-deterministische manuelle Eingaben zurückging -- abweichende Schreibweisen von Personen- und Ortsnamen, uneinheitliche Datumsformate und Mehrfacheinträge mit leichten Varianten. Eine systematische Übersicht über die Bereinigungsschritte findet sich bei \cite{castorp2024-cleanup}.

Im zweiten Schritt wurde gemeinsam mit den fachlichen Verantwortlichen ein Zielmodell auf Basis von CIDOC CRM entwickelt. Die Modellierung wurde bewusst minimal gehalten, um die anschließende Anreicherung zu erleichtern. Drittens erfolgte die eigentliche Transformation: Aus den profilierten Quelldaten wurden über ein in Python implementiertes ETL-Pipeline-Tool RDF-Tripel erzeugt, die direkt in den von WissKI verwendeten GraphDB-Triplestore eingespielt werden konnten. Der vierte Schritt -- die Qualitätssicherung -- erfolgte mit SHACL-basierten Validierungsregeln und einer stichprobenartigen Re-Erfassung durch das Fachteam.

\section{Werkzeugkette und Daten-Mapping}
Die folgende Tabelle gibt einen Überblick über die im Projekt eingesetzten Werkzeuge und ihre jeweilige Rolle.

\begin{table}[ht]
    \centering
    \begin{tabular}{lll}
        \toprule
        \textbf{Phase} & \textbf{Werkzeug} & \textbf{Aufgabe} \\
        \midrule
        Profilierung & OpenRefine    & Inkonsistenzen, Cluster, Statistiken \\
        Modellierung & Protégé       & CRM-Profilentwurf \\
        ETL          & Python / RDFLib & Transformation, Tripelerzeugung \\
        Validierung  & PySHACL        & Schemaregeln, Berichte \\
        Persistenz   & GraphDB        & Triplestore-Backend für WissKI \\
        Frontend     & WissKI / Drupal & Eingabe und Recherche \\
        \bottomrule
    \end{tabular}
    \caption{Werkzeuge in der Migrationspipeline.}
    \label{tab:p02-werkzeuge}
\end{table}

Eine zentrale Designentscheidung war, jeden Migrationsschritt deterministisch und reproduzierbar zu gestalten -- ETL-Skripte, Konfigurationsdateien und SHACL-Shapes liegen versioniert im Projektrepositorium vor, sodass eine Re-Migration mit aktualisierten Quelldaten ohne manuelle Eingriffe möglich ist.

\section{Ergebnisse}
Die Migration konnte innerhalb von vierzehn Monaten abgeschlossen werden. 11\,840 von 12\,400 ursprünglichen Datensätzen wurden vollständig automatisch überführt; die verbleibenden 560 mussten manuell nachbearbeitet werden -- in der Regel wegen widersprüchlicher Aussagen über dieselbe Entität in unterschiedlichen Quellen. Die anschließende fachliche Anreicherung ist deutlich schneller geworden, weil die Datenbasis nun in einem konsistenten Modell vorliegt und über eine SPARQL-Schnittstelle abfragbar ist; vgl.\ \cite{drechsler2025-erfahrungsbericht}.

\section{Fazit}
Die Erfahrung des Projekts lässt sich in einer einfachen Faustregel zusammenfassen: \emph{Profilieren, bevor man modelliert; modellieren, bevor man transformiert; validieren, bevor man veröffentlicht.} Die Konsequente Anwendung dieser Reihenfolge -- bei aller Versuchung, schon vorab in das Triple-Store-System hineinzuoptimieren -- hat sich als entscheidender Faktor für den Projekterfolg erwiesen.

\end{beitrag}
